\scanpage{104}
\mspage{\handwritten{$6'$}}

In a more suggestive way, we could say \deleted{\unsure{that}}\unsure{from|thus} this experimenting that the family of discs $(D_i)_{i\in\underline{N}}$, together with the maps $s,t$ and a lot of extra structure which enters into the picture step by step, is what we would like to call a \uline{co}-$\infty$-\uline{groupoid in the category} (Top) of \unsure{topologica|topological} spaces (namely an $\infty$-groupoid in the dual category $(\mathrm{Top})^o$), and that the structure of $\infty$-category on $F_*(X)$ we want to describe is the transform of this co-structure into an $\infty$-groupoid, by the contravariant functor from (\unsure{Tp|Top}) to (Sets) defined by $X$. The \inserted{\handwritten{(iterated)}} amalgamated sums in (Top) which allow to glue together \deleted{\unsure{the|these}} various $D_i$'s using the $s$ and $t$ maps between them, namely the corresponding fibered products in $(\mathrm{Top})^o$, are indeed transformed by the functor $h_X$ into fibered products of (Sets).
The suggestion is, moreover, that if we view our co-structure in (Top) as a co-structure in the subcategory of Top, say $B_\infty$, whose objects are the \deleted{\unsure{$D_i$}}cells $D_i$ and the amalgamated sums \deleted{\unsure{r|s}}built up with these which step-wise enter into play, and whose arrows are all those arrows which are introduced stepwise to define the co-structure, and all compositions of these - that \deleted{\unsure{it}}\handwritten{this} \deleted{\unsure{t}}should be \deleted{\illegible{1}}\uline{\handwritten{the} universal structure} in the sense dual to the one we have been contemplating before; or what amounts to \deleted{th} the same, that the corresponding $\infty$-groupoid structure in the dual category $B_\infty^o$ is ``universal'' - which means essentially that it is none other than the universal structure in the category $\underline{C}_\infty$ we are after.

Whether or not this expectation \deleted{\unsure{be|is}}\inserted{\handwritten{will turn out to be}} correct (I believe it is\handwritten{$^{*}$}), we should
% slanted
\margin{\handwritten{* (23.2) But no longer now - and \unsure{I} do not really care!}}
be aware that, while the successive introduction of maps between the cells $D_i$ and their ``standart'' amalgamated sums (which we will define precisely below) depends at every stage on arbitrary choices, the catego\deleted{ry}ries $\underline{C}_n$ and their limit $\underline{C}_\infty$ do not depend on any \inserted{of these} choices ; assuming the expectation is correct, this means that up to (unique) isomorphism, the category $B_\infty$ (and each of the category $B_n$ of which it appears as the direct limit) is independant of those choices - the isomorphism between two such categories transforming any one choice made for the first, into the corresponding choice made for the second. Also, while this expectation was of course the crucial motivation leading to the explicit description of $\underline{C}_0$ and of the \inserted{inductive} step from $\underline{C}_n$ to $\underline{C}_{n+1}$, this description seems to me a reasonable one and in any case it makes a formal sense, \inserted{\handwritten{quite}} independently of whether the expectation proves valid or not

\handwritten{7.} Before pursuing, it is time to give a more complete discription of the primitive structure on $(D_i)$, as embodied by the maps $s,t$, which I will now denote by
\[
{\unsure{\phi|\delta|\varnothing}}_i^+,\,\phi_i^-:D_i\rightrightarrows D_{i+1}.
\]
It appears that these maps are injective, that their images make up
% archive sheet number: 103
